%
%
%
%
\section{Open Problems}\label{sec:open}
\begin{enumerate}
\item \emph{One sequence certified at both kinds of levels.} No digital sequence is known to us that is certified both at every $m\equiv-1\pmod s$ with the constants of Theorem~\ref{thm:residue} and at every full level with $1-O(1/p)$. For exact designed relations of bounded depth next to exact skeleton relations this is excluded (Corollary~\ref{cor:nodesign}), and more generally exact relations of bounded depth leave room for about one relation per level (Theorem~\ref{thm:count}); for windows chosen together with a corrected skeleton relation in a self-similar scramble, Proposition~\ref{prop:joint} excludes it in the scope stated there (product windows, covering of every line, the correction window of Theorem~\ref{thm:corrected}). Open are selection rules that look at later digits, windows that are not products or that use several digits, certificates other than the one-digit relaxation, scrambles that are not self-similar, and digital sequences that are not scrambles of a frame of finite places.
\item \emph{Every full level with constants tending to $1$ at all primes.} At every full level the constants $1-O(1/p)$ are certified only from large primes on: $1-17/p$ for $p>3677$ in dimension two (Corollary~\ref{cor:HB17}), and $1-C''_s/p$ from $p_0(s)$ on in dimension $s\ge3$, with $p_0(3)\le1.97\cdot10^7$ and $p_0(4)\le1.44\cdot10^{16}$ (Theorems~\ref{thm:A} and~\ref{thm:A3}), because one hierarchy of step sizes serves all columns. The block-periodic residue sequences reach $1-2/p$ from $p=3$ ($s=2$) and $p=5$ ($s=3$) on, but only at every $s$-th full level, and lose one $1/p$ against the $1-1/p$ of the nets of \cite{Gru26geo}%
, as the deeper digits of a coordinate depend on its own cell. Possible routes: windows on the free multipliers of block-periodic scrambles with period $\ell>s$, to certify the full levels between the design levels; coverings of lattice type for the staggered windows, to lower $p_0(s)$.
\item \emph{Base five in dimension three at every full level.} Base $5$ is covered only on the level family $9\N$ (Theorem~\ref{thm:block}(b)); Remark~\ref{rem:base5} lists what was tried for every full level: symmetric frames have no three-dimensional instance over $\F_5$, the staggered residue sequences have explicit close pairs at deep levels, and several one- and two-digit certificate classes were reported unsatisfiable by a SAT solver in the instances solved (evidence about those classes only). General two-digit windows, history-dependent bounds with several digits (as in Appendix~\ref{app:p7}), and non-symmetric algebraic frames remain unexplored.
\item \emph{Random linear scrambles of Sobol'.} Every Sobol' sequence fails (Corollary~\ref{cor:sobol}), and so does every scramble that leaves one coordinate a radical inverse up to a scalar (Corollary~\ref{cor:unscrambled}); a fixed two-sided linear scramble of the two-dimensional Sobol' sequence is quasi-uniform (Theorem~\ref{thm:B2}). Whether a random linear matrix scramble of the Sobol' sequence is quasi-uniform, almost surely or with positive probability, is open: the bound $O_P(\log N)$ for the mesh ratio of randomly linearly scrambled binary nets that satisfy a uniform coincidence condition, as the Sobol' prefixes do \cite[Thm.~5.5, Lem.~4.10]{Suz26}, does not decide it. How fast the separation of the Sobol' and Niederreiter sequences decays in dimension $s\ge3$ is open as well: Section~\ref{sec:fail} gives explicit close pairs along one sequence of sizes, and the rate $N^{-1/(s-1)}$ only for the Faure sequence (Remark~\ref{rem:faurerate}).
\item \emph{The limit of the residue-sequence constant.} The best certified constant $\max_vp^{-1/s}\ljit(\Lambda_v)$ of the residue sequences (Table~\ref{tab:residue}) appears to tend to $\sqrt{\gamma_s}$, the square root of the Hermite constant; this is the open limit of the skeleton nets of \cite{Gru26geo}%
, and an answer there transfers verbatim. The upper bound $\ljit(\Lambda_v)\le\sqrt{\gamma_s}\,p^{1/s}$ is immediate, since $\Lambda_v$ has index $p$. A route for the lower bound that we have not verified: the lattices $\Lambda_v$ are the sublattices of index $p$ of $\Z^s$ (those whose check vector has no zero entry), and an equidistribution statement for these sublattices, rescaled to covolume $1$, in the space of lattices as $p\to\infty$ could give the limit, since the jitter correction costs at most $\sqrt s\,p^{-1/s}$; we have neither checked whether such a statement is available in the form needed nor worked out the details.
\item \emph{Implied skeleton relations in dimension $s\ge3$.} In dimension two a design can imply a skeleton relation at the level after the design level (Proposition~\ref{prop:implied}, Corollary~\ref{cor:gauss}), and with the skeleton multipliers of Theorem~\ref{thm:blocktwo} the two families meet the density bound of Corollary~\ref{cor:density} (Remark~\ref{rem:gaussfrontier}). For $s\ge3$ the conditions of Proposition~\ref{prop:implied} can be met algebraically by the free diagonal, but the remark after it forces an entry of order $p^{1/s}$ in every nonzero column of the design for a good skeleton constant, and then the design constant is no longer $1-O(1/p)$. No example for $s\ge3$ has been constructed; a search over places, check vectors and crosses has not been run.
\item \emph{The true full-level separation of the staggered residue sequence.} The certified bound $(\lfloor p/2^{s-1}\rfloor-1)/p\approx2^{1-s}$ of Corollary~\ref{cor:dyadic} is a property of the one-coordinate certificate, not of the sequence: for $s=4$ and $s=5$ the exact separations of the first full level $P_s$ in the cases computed lie between $0.1985$ and $0.4837$ (Section~\ref{sec:brickwall}). Whether $\inf_K\sep{P_{sK}}$ stays of order $1$ for $s\ge4$, or decays to the order $2^{1-s}$ of the bound, is open; a certificate that uses all coordinates would be needed.
\item \emph{Small primes and prime powers.} The primes $s\le p<2^s$ in dimension $s\ge4$ are not covered at every full level (on $s^2\N$, Theorem~\ref{thm:block} covers them from $p=7$, $11$, $17$ on for $s=4,5,6$), and whether residue sequences are quasi-uniform there is not addressed. The staggered residue sequences keep separations well away from $0$ at the first levels below $2^s$ (for instance $0.48$ to $0.52$ for $K\le7$ at $s=2$, $p=3$; exact), but the close pairs found at deep levels for $s=3$, $p=5$ (Remark~\ref{rem:base5}) show how little such data say. Prime powers beyond Proposition~\ref{prop:primepower}, the residue sequence over $\F_q$ (the proof of Theorem~\ref{thm:residue} uses that $q$ is prime), and the more general Halton-type sequences of Hofer in rational bases \cite{Hof18}, to which \cite{GHS26} refer (outside the definition treated in Corollary~\ref{cor:halton}), are open.
\item \emph{Quality at reachable sizes, integration, and base two.} The certified bounds are loose at small primes, the Faure sequence is better at its first full level, and in the scans of Section~\ref{sec:practice} the choice of places, scale and check vector moved $\mathcal J$ more than the certified constant did. Criteria that target the first levels while keeping a certified bound, and a measured study of integration error, are open. A $(0,s)$-sequence in base $2$ needs $s\le2$; the binary sequences of Theorems~\ref{thm:K4} and~\ref{thm:S4} are quasi-uniform in dimensions three and four, with certified bounds at the sparse levels $N=2^{6K}$ and $2^{8K}$ and small constants. Binary sequences with certified bounds at denser levels, or with constants comparable to those of large prime bases, are open.
\end{enumerate}
